%% kapa: exam=1 round=21 year=2010 period=2 form=A subject=accounting subject_ko=회계학 seq=27 number=27
%% source: Q-Net 감정평가사 기출문제 (cst00309) artlSeq=5024181 file=extracted/1008/5024181/2010년 제21회 감정평가사 1차시험 기출문제/2교시.hwp
%% conversion: kordoc 4.24.51 -> tools/gen1.py (auto); review=proofread-gpt-sol; flags=none
\begin{question}{27}
(주)한국은 20×9년 중 외국에 있는 금융기관으로부터 만기 3년의 외화표시 부채 ＄1,000을 차입하여 상각후원가로 평가하였다. 차입 일자에 달러 현물의 마감환율은 ＄1 당 ￦1,000이었다. 20×9년 말 현재 달러의 외화현물에 대한 마감환율이 ＄1 당 ￦1,100으로 상승하였다면, 장기외화차입금에 대한 환율의 상승분 ￦100,000은 포괄손익계산서에 어떻게 보고되는가?
\begin{choices}
  \item 당기이익으로 보고된다.
  \item 기타포괄이익으로 보고된다.
  \item 공정가치평가손실로 보고된다.
  \item 당기손실로 보고된다.
  \item 기타포괄손실로 보고된다.
\end{choices}
\end{question}
\answer{4}
